\documentclass[11pt]{letter}
\usepackage[margin=1in]{geometry}

\begin{document}

\begin{letter}{}

\opening{Dear Members of the Search Committee,}

I am writing to apply for a tenure-track faculty position at Johns Hopkins University’s new School of Government and Policy. I am a Ph.D. candidate in Economics at Johns Hopkins University. My research lies at the intersection of macroeconomics, public finance, and inequality, with a focus on combining structural modeling with micro and panel data to study distributional outcomes and the design of economic policy.

My job market paper develops and estimates a heterogeneous-agent macroeconomic model with differences in returns across households to study wealth inequality and the welfare effects of tax policy. Using micro-level data from the Survey of Consumer Finances, I estimate return heterogeneity to match observed wealth dynamics and simulate counterfactual policy reforms. A central finding is that the interaction between income risk and return heterogeneity generates realistic wealth dispersion and marginal propensities to consume, allowing the model to match both cross-sectional inequality and aggregate responses to policy. 

More broadly, my research program integrates quantitative macroeconomic modeling with empirical analysis using detailed microdata. In ongoing work using the Health and Retirement Study, I employ panel data methods and causal identification strategies to study the relationship between trust, financial behavior, and household returns. This work connects individual-level behavior to broader questions about inequality, financial decision-making, and economic resilience.

My research aligns closely with several of the School’s core areas of interest, particularly Cities and Communities and Governance. I am especially interested in how local economic conditions, labor market dynamics, and financial constraints shape household outcomes and interact with policy design. In addition, I am developing work on the distributional effects of macroeconomic shocks, including a project examining how oil price shocks transmit to income, consumption inequality, and labor market outcomes using micro-level data. These projects contribute to a broader agenda focused on understanding how policy and economic environments shape inequality and economic opportunity.

The School of Government and Policy’s emphasis on interdisciplinary collaboration, policy engagement, and intellectual pluralism makes it an ideal environment for my work. I am particularly excited about the opportunity to contribute to a new institution designed to bridge rigorous research with real-world policy impact, and to engage with scholars across disciplines working on questions related to governance, economic development, and technological change.

I also have experience teaching Principles and Intermediate Macroeconomics, and I am comfortable teaching both undergraduate and graduate courses. My teaching emphasizes building economic intuition alongside quantitative tools, helping students connect theory, data, and policy in a clear and accessible way.

Thank you for your consideration. I would be thrilled to contribute to the School’s research, teaching, and policy mission.

\closing{Sincerely,}

\vspace{0.5cm}

Decory Edwards

\end{letter}

\end{document}